\documentclass[12pt]{article}
\usepackage[margin=1in]{geometry}
\usepackage{amsmath,amssymb,amsthm,booktabs,natbib,microtype}
\newtheorem{theorem}{Theorem}
\newtheorem{proposition}[theorem]{Proposition}
\newtheorem{corollary}[theorem]{Corollary}
\newtheorem{remark}[theorem]{Remark}
\newcommand{\norm}[1]{\lVert#1\rVert}
\title{The Midgap Circle for Normal Packet Blocks\\\large An Exact Gap--Feedback Criterion and Its Nonnormal Boundary}
\author{Prime Dynamics Theory Program}
\date{July 2026}

\begin{document}
\maketitle
\begin{abstract}
The Schur theorem of RH-163 reduces packet-to-Riesz rank preservation to
$adbc<1$, but leaves the contour and two resolvent bounds unspecified.  We
solve this geometric problem for a benchmark class.  Suppose the packet and
complement blocks are normal, the packet spectrum lies in
$|z-\mu|\le\rho$, and the complementary spectrum lies in
$|z-\mu|\ge R>\rho$.  Among centered circles, the radius
\[
 s_*=(\rho+R)/2
\]
minimizes the Schur feedback.  The optimized rank condition is exactly
\[
 2\sqrt{bc}<R-\rho.
\]
We also give the resulting projector envelope and prove sharpness of the
midpoint for the available disk data.  A 512-system finite audit has zero
rank-count and resolvent-bound failures.  The theorem deliberately stops at
normal blocks: for a nonnormal operator, spectral distance alone cannot
replace a resolvent estimate.  No physical normality or all-level
pseudospectral theorem is claimed.
\end{abstract}

\section{Disk-separated normal blocks}
Let
\[
 A=\begin{pmatrix}A_P&B\\C&A_Q\end{pmatrix}
\]
on an orthogonal packet decomposition.  Assume $A_P$ and $A_Q$ are normal
and, for a center $\mu\in\mathbb C$,
\[
 \sigma(A_P)\subset\{|z-\mu|\le\rho\},\qquad
 \sigma(A_Q)\subset\{|z-\mu|\ge R\},\quad R>\rho.
\]
Put $b=\norm B$ and $c=\norm C$.  For $\rho<s<R$, let
$\Gamma_s=\{|z-\mu|=s\}$.

Normality gives the exact spectral-distance resolvent formula.  Therefore
\[
 a(s):=\sup_{\Gamma_s}\norm{(z-A_P)^{-1}}\le\frac1{s-\rho},
\quad
 d(s):=\sup_{\Gamma_s}\norm{(z-A_Q)^{-1}}\le\frac1{R-s}.
\]

\section{Optimal midpoint theorem}
\begin{theorem}[Midgap feedback optimum]\label{thm:midgap}
Among centered circles $\Gamma_s$, the upper bound
\[
 \kappa(s)=\frac{bc}{(s-\rho)(R-s)}
\]
is minimized uniquely at
\[
 s_*=(\rho+R)/2,
\]
where
\[
 \kappa_* = \frac{4bc}{(R-\rho)^2}.
\]
Consequently $\Gamma_{s_*}$ certifies packet-rank preservation whenever
\[
 \boxed{\ 2\sqrt{bc}<R-\rho.\ }
\]
\end{theorem}
\begin{proof}
The denominator is
\[
 (s-\rho)(R-s)
 =\frac{(R-\rho)^2}{4}
  -\left(s-\frac{R+\rho}{2}\right)^2.
\]
It is uniquely maximized at the midpoint.  Substitution gives $\kappa_*$.
If $\kappa_*<1$, the RH-163 Schur homotopy theorem preserves the enclosed
rank \citep{WangRH163}.
\end{proof}

The constant is sharp for this disk information.  Scalar packet and
complement blocks placed at the two nearest radial boundary points attain
both resolvent denominators.  At $4bc=(R-\rho)^2$, a two-mode matrix can have
an eigenvalue on the midpoint circle.

\section{Projector envelope}
Let $g=R-\rho$, $s_*=(R+\rho)/2$, and
$\kappa=4bc/g^2<1$.  At the midpoint, $a=d=2/g$.  The scalar block estimate
of RH-163 yields
\[
 \norm{\Pi-P}\le
 \frac{s_*}{1-\kappa}
 \left\|
 \begin{pmatrix}
 2\kappa/g&4b/g^2\\
 4c/g^2&2\kappa/g
 \end{pmatrix}\right\|_2.
\]
This graph bound depends on the absolute contour radius $s_*$ as well as the
gap.  Moving both spectral sets far from the origin leaves the rank gate
unchanged but lengthens the circle; a noncircular contour or a directional
graph estimate may then be preferable.

\section{Why the theorem stops at normality}
For a normal operator,
$\norm{(z-A)^{-1}}=\operatorname{dist}(z,\sigma(A))^{-1}$.
For a nonnormal matrix, the resolvent can be arbitrarily larger while the
spectrum and radial gap are fixed \citep{TrefethenEmbree2005}.  A Jordan
block already exhibits polynomial resolvent amplification.  Therefore the
condition $2\sqrt{bc}<g$ is not a physical prime-dynamics theorem unless
normality, similarity conditioning, or an independent pseudospectral
envelope is proved.

One can retain the geometry without normality by replacing $\rho$ and $R$
with certified pseudospectral barriers that directly imply the displayed
resolvent bounds.  That replacement is an input, not a consequence of the
eigenvalue plot.

\section{Conditioned pseudospectral inflation}
There is a useful intermediate statement between normality and a completely
free resolvent.  Suppose that on every centered circle one has certified
envelopes
\[
 \sup_{\Gamma_s}\norm{(z-A_P)^{-1}}
 \le\frac{K_P}{s-\rho},\qquad
 \sup_{\Gamma_s}\norm{(z-A_Q)^{-1}}
 \le\frac{K_Q}{R-s},
\]
where $K_P,K_Q\ge1$ are fixed conditioning factors.

\begin{corollary}[Inflated midgap gate]
The midpoint circle remains optimal for these envelopes and certifies rank
whenever
\[
 2\sqrt{K_PK_Qbc}<R-\rho.
\]
\end{corollary}
\begin{proof}
The Schur product is bounded by
$K_PK_Qbc/[(s-\rho)(R-s)]$.  The same quadratic denominator as in
Theorem~\ref{thm:midgap} is maximized at the midpoint, where the product is
$4K_PK_Qbc/(R-\rho)^2$.
\end{proof}

This corollary quantifies the price of nonnormality when a uniform similarity
or pseudospectral constant is available.  If $K_PK_Q$ grows with scale, the
geometric gap may widen while the usable margin still collapses; both factors
must therefore be recorded in the physical ledger.

\section{Audit and route consequence}
The audit samples 512 diagonal normal block pairs with complex spectra in
the two radial regions.  Directed couplings are scaled so the midpoint
feedback is below $0.55$.  On 1024 contour nodes it verifies both resolvent
bounds, and direct diagonalization confirms exactly the packet rank inside
the circle.  No failure occurs.  Floating-point rank counts are illustrative
regression checks, not interval proofs.

RH-165 removes contour optimization from the normal benchmark and supplies a
sharp geometric falsifier.  The physical family remains nonnormal, so the
next steps must obtain directed Ritz residuals and then certify continuous
resolvent envelopes.  The paper does not close physical R, Gate A, or any
later Hilbert--Polya interface.

\bibliographystyle{plainnat}
\bibliography{references}
\end{document}
